% 《大学物理学学习指导》第5版自测题1：计算1.3.1—1.3.5，改错1.4.1。
% 原书印刷47—48／PDF物理55—56页；图示按原书方向转述。
\subsection{自测题1：动量、功与机械能计算}

\sourcetest{1.3.1}{47}{55}
\problemtext 质量$m$的小块从高度$h$的光滑曲面滑下，
曲面固定在质量$M$的可水平自由滑动楔块上。
楔面从左侧$x=0$延伸至右端$x=L$，
初始系统静止，求滑块落至底部期间楔块水平移动距离。
\workings
取楔块相对地面位移$\Delta X$（向右正），
小块的绝对水平位移是$\Delta X+L$。
全程没有外界水平力，系统质心的水平位置不变，
\[
M\Delta X+m(\Delta X+L)=0.
\]
因此
\[
\boxed{\Delta X=-\frac{m}{M+m}L}.
\]
楔块向左移动，距离大小为$mL/(M+m)$。
\crosscheck 位移由水平质心守恒给出，
\emph{与}曲面下落高度$h$的具体数值无关；
$h$决定速度与能量，不决定上述质心位移结果。

\sourcetest{1.3.2}{47}{55}
\problemtext 两艘质量同为$M$的船初始在静水中静止，
质量$m$的人以对地水平速度$v$从第一艘跳到第二艘，
忽略水阻，求两艘船在跳离及落入后各自速度。
\workings
以跳跃方向为正。
人刚离开船1时，船1与人的总水平动量守恒：
$MV_1+mv=0$，故
\[
\boxed{V_1=-\frac mM v}.
\]
人登上最初静止的船2并与它共同运动时，
对人+船2的瞬时完全非弹性相互作用应用动量守恒，
$(M+m)V_2=mv$，故
\[
\boxed{V_2=\frac{m}{M+m}v}.
\]
\crosscheck 船1向与人相反方向退行，船2沿人跳的方向前进。
人给定的$v$是\emph{相对于地面}，不能误当相对船的速度。

\sourcetest{1.3.3}{47--48}{55--56}
\problemtext 皮带以恒速$v$沿水平方向运行，
无初始水平速度的砂从静止漏斗落到带上。
单位时间落砂质量$\dot m=dM/dt$，
求维持皮带速度所需水平拉力和电机输出机械功率；
再取$\dot m=20\,\mathrm{kg/s},v=1.5\,\mathrm{m/s}$。
\workings
在$dt$内有$dm=\dot m\,dt$的砂从水平速度0
加速到皮带速度$v$，动量增量为$v\,dm$，
所以皮带必须施于砂的平均水平力
$F=v\,dm/dt=\dot m v$，
相应电机对皮带的机械功率
\[
\boxed{P=Fv=\dot m v^2}.
\]
数值为
\[
\boxed{F=20(1.5)=30\,\mathrm N},\qquad
\boxed{P=20(1.5)^2=45\,\mathrm W}.
\]
\crosscheck 砂获得的平动动能率仅为
$\tfrac12\dot m v^2=22.5$W；
其余机械功率在砂与带相对滑动时转为热。
因此\emph{电机机械输出功率}不是简单等于砂动能增长率。

\sourcetest{1.3.4}{48}{56}
\problemtext 质量$M$的人手持质量$m$物体，
以初速$v_0$、仰角$\alpha$从地面向前跳跃，
到弹道最高点时将物体以相对人向后的水平速度$u$抛出。
求抛出物体相对不抛时，此人的水平落地点增加距离。
\workings
跳跃全程竖直初速度$v_0\sin\alpha$，
到最高点历时$t_h=v_0\sin\alpha/g$。
抛出前人和物体的水平速度均为$v_0\cos\alpha$。
水平动量在抛出的短暂时间内守恒；
设抛后人速度$v_M$，物体速度为$v_M-u$，
\[
(M+m)v_0\cos\alpha=Mv_M+m(v_M-u).
\]
解得人新增水平速度
\[
\Delta v_M=v_M-v_0\cos\alpha=\frac{mu}{M+m}.
\]
从最高点到落地仍需$t_h$，
故人与\emph{不抛物体}时落地点的差为
\[
\boxed{\Delta x=\frac{muv_0\sin\alpha}{(M+m)g}}.
\]
\crosscheck 题中的$u$是\emph{相对于人}的抛速。
若改成物体对地的抛速，水平动量方程应相应修改。

\sourcetest{1.3.5}{48}{56}
\problemtext 质量$M$的空心物体A在光滑平面上，
内部水平弹簧劲度系数$k$连接内壁，
质量$m$的物体B紧靠弹簧但未固定。
初始弹簧压缩$\Delta x$，全系统由静止释放。
当B与弹簧分开后，求A的反冲速度大小。
\workings
两物体沿水平方向无外力，
分离后动量大小相等而方向相反，
设A反冲速度大小$u$、B速度大小$v$：
\[
Mu=mv.
\]
无摩擦，弹簧初始势能在恢复原长时化为两物体动能，
\[
\frac12k(\Delta x)^2
=\frac12Mu^2+\frac12mv^2.
\]
将$v=Mu/m$代入求得
\[
\boxed{u=\Delta x\sqrt{\frac{km}{M(M+m)}}}.
\]
\crosscheck 极限$M\to\infty$时$u\to0$；
与动量守恒一致。须把A的整个质量作为反冲体质量，
不能只用内壁弹簧质量。

\subsection{自测题1：改错题}

\sourcetest{1.4.1}{48}{56}
\problemtext 轻弹簧原长时挂一质量为$m$的物体，
从静止释放后向下运动。
某解答认为“合力为零时，重力势能减少等于弹性势能增加”，
由$mg y_0=\tfrac12ky_0^2$求出$y_0=2mg/k$，
但又把$y_0$当成\emph{合力为零}的伸长量。指出并改正。
\workings
以向下伸长量$y$为正，
系统的加速度为$a=g-ky/m$，
所以\emph{合力为零}时
\[
\boxed{y_{\rm eq}=\frac{mg}{k}},
\]
该点是静态平衡点，却不是从原长释放过程的
最大下降点；到这里时速度尚不为零。
机械能守恒（初始$y=0,v=0$）给
\[
\frac12mv^2+\frac12ky^2-mgy=0.
\]
在平衡点代入$y=mg/k$得
\[
\boxed{\frac12mv_{\rm eq}^2=\frac{m^2g^2}{2k}>0}.
\]
而下降的\emph{最低点}满足$v=0$，
由$mg y_{\max}=\frac12ky_{\max}^2$解非零根，
\[
\boxed{y_{\max}=\frac{2mg}{k}}.
\]
\crosscheck 数值$2mg/k$对应\emph{速度为零的最低点}，
并非\emph{合力为零的平衡点}；
原解的错误是将这两个不同事件混为一谈。
